\chapter{Fundamentação Teórica}
\label{cap:fundamentacao}

\abbrev{MPI-IO}{MPI Input/Output}
\abbrev{OST}{Object Storage Target}
\abbrev{MDS}{Metadata Server}
\abbrev{MDT}{Metadata Target}
\abbrev{OSS}{Object Storage Server}
\abbrev{RDMA}{Remote Direct Memory Access}
\abbrev{AWS}{Amazon Web Services}
\abbrev{FSx}{Amazon FSx for Lustre}
\abbrev{EFA}{Elastic Fabric Adapter}
\abbrev{ONS}{Operador Nacional do Sistema Elétrico}
\abbrev{PMO}{Programa Mensal da Operação}
\abbrev{SIN}{Sistema Interligado Nacional}
\abbrev{LNCC}{Laboratório Nacional de Computação Científica}
\abbrev{SINAPAD}{Sistema Nacional de Processamento de Alto Desempenho}

\section{MPI e protocolos de comunicação}
O \textit{Message Passing Interface} (MPI) é o padrão para programação
paralela baseada em troca de dados entre processos, sendo amplamente utilizado em
\textit{clusters} e supercomputadores. Um aspecto essencial da sua
implementação são os protocolos de comunicação: o \textit{eager}, que envia
pequenos blocos de dados de forma imediata, sem necessidade de confirmação do
receptor, favorecendo baixa latência; e o \textit{rendezvous}, que aguarda o
\textit{handshake} do receptor antes de transmitir blocos de dados maiores,
reduzindo \textit{overhead} de memória e aumentando a previsibilidade do
desempenho em cargas intensivas~\cite{mpiCite}. A escolha entre os
protocolos depende fortemente do tamanho do bloco de dados e da disponibilidade de
\textit{buffers}: enquanto o \textit{eager} utiliza \textit{buffers}
internos para armazenar temporariamente dados até que o receptor os
processe, o \textit{rendezvous} elimina essa necessidade para blocos de dados
grandes, evitando consumo excessivo de memória e garantindo maior
estabilidade em aplicações de larga escala~\cite{Wea:99,Rrw:05}.

\section{MPI-IO e operações paralelas de arquivos}
O MPI também estende suas funcionalidades à entrada e saída de dados por
meio da especificação MPI-IO~\cite{mpiioCite}, que define uma API
padronizada para operações paralelas de arquivos. Essa interface possibilita
que múltiplos processos acessem simultaneamente grandes volumes de dados,
permitindo otimizações como \textit{collective I/O}, técnicas de agregação e
\textit{data sieving}, que reduzem o número de operações pequenas e melhoram
a eficiência do uso do sistema de arquivos. Entre as chamadas mais
utilizadas, destacam-se \texttt{MPI\_File\_write\_at}, que realiza uma
escrita independente em um deslocamento específico do arquivo, permitindo
controle preciso sobre onde os dados serão gravados, e
\texttt{MPI\_File\_write\_all}, que executa uma operação coletiva em que
todos os processos participantes escrevem de forma coordenada. Enquanto a
primeira é adequada para acessos irregulares ou não sincronizados, a segunda
explora a comunicação coletiva para reduzir \textit{overhead} de metadados e
alinhar requisições, sendo mais eficiente em ambientes com sistemas de
arquivos distribuídos como o Lustre.

\section{Sistema de arquivos paralelo Lustre}
Projetado especificamente para cargas massivas de HPC, ele distribui os
dados em \textit{Object Storage Targets} (OSTs) e os metadados em
\textit{Metadata Servers} (MDS), permitindo acesso paralelo e balanceado a
múltiplos clientes. Além disso, o Lustre utiliza a técnica de
\textit{striping}, na qual um arquivo é dividido em blocos
(\textit{chunks} ou \textit{stripes}) distribuídos entre diferentes OSTs.
O tamanho do \textit{stripe} e o número de OSTs utilizados podem ser
configurados de acordo com o padrão de acesso da aplicação, influenciando
diretamente no \textit{throughput} e na escalabilidade: \textit{stripes}
maiores tendem a favorecer grandes operações sequenciais, enquanto
\textit{stripes} menores podem beneficiar acessos concorrentes e
distribuídos~\cite{Pjb:02}. Graças a essa flexibilidade e alto desempenho,
o Lustre tornou-se amplamente adotado em supercomputadores listados no
TOP500, constituindo um elemento-chave para aplicações científicas e
industriais que exigem desempenho extremo em E/S.

\section{Adaptadores de rede em HPC: Ethernet e InfiniBand}
Os adaptadores de rede são componentes críticos para o desempenho em
sistemas de computação de alto desempenho, sendo a Ethernet e o InfiniBand
as tecnologias mais utilizadas. A Ethernet, tradicionalmente projetada para
redes locais corporativas, evoluiu para padrões de alta velocidade, como
40/100/400~Gbps, oferecendo elevada taxa de transferência
(\textit{bandwidth}), embora apresente latências relativamente maiores,
frequentemente na ordem de dezenas de microssegundos. Já o InfiniBand foi
desenvolvido especificamente para ambientes de HPC, combinando altas
larguras de banda --- que podem ultrapassar 400~Gbps em implementações
modernas --- com latências extremamente baixas, tipicamente abaixo de
2~$\mu$s, graças a recursos como comunicação RDMA (\textit{Remote Direct
Memory Access}). Essa diferença torna o InfiniBand preferido em
supercomputadores e \textit{clusters} científicos, onde a baixa latência na
troca de dados entre processos é tão relevante quanto a banda disponível, enquanto a
Ethernet, pelo menor custo e ampla compatibilidade, mantém espaço em
sistemas híbridos e \textit{datacenters}~\cite{Fds:03,RlGmmts:05}.

\section{Computação em nuvem para HPC: AWS}
A Amazon Web Services (AWS) tem se consolidado como uma das principais
plataformas de computação em nuvem para aplicações de \textit{High
Performance Computing} (HPC), oferecendo recursos sob demanda que permitem
a execução de cargas de trabalho intensivas em computação e E/S. Entre os
serviços disponíveis, destacam-se as instâncias otimizadas para HPC, como
as famílias C, H e P, que fornecem alto poder de processamento aliado a
interconexões de baixa latência por meio do \textit{Elastic Fabric Adapter}
(EFA), permitindo a utilização de bibliotecas como MPI em escala. Para
atender aplicações intensivas em dados, a AWS integra sistemas de
armazenamento de alto desempenho, como o Amazon FSx for Lustre, que
viabiliza acesso paralelo e em larga escala, semelhante a sistemas de
arquivos utilizados em supercomputadores tradicionais. Essa combinação de
infraestrutura elástica, rede de alto desempenho e suporte a ferramentas
amplamente usadas em HPC torna a nuvem da AWS uma alternativa competitiva
frente a \textit{clusters} locais, especialmente em cenários que demandam
escalabilidade e flexibilidade~\cite{Krlkssj:10,Hmns:20}.

\section{O Programa Mensal da Operação (PMO)}
No setor elétrico brasileiro, o Programa Mensal da Operação (PMO) é um
documento técnico elaborado pelo Operador Nacional do Sistema Elétrico
(ONS) que estabelece as diretrizes de curto prazo para a operação do
Sistema Interligado Nacional (SIN). O PMO contempla projeções de carga,
disponibilidade de geração, restrições de transmissão e condições
hidrológicas, servindo como referência para a programação da operação
hidrotérmica de forma a garantir o equilíbrio entre oferta e demanda de
energia elétrica. Além disso, define metas de despacho de usinas
hidrelétricas e termelétricas, contribuindo para a segurança energética e
para a otimização do uso dos recursos disponíveis. Por sua relevância, o
PMO orienta não apenas as decisões operativas do ONS, mas também agentes
do setor, incluindo geradores, comercializadores e distribuidores, sendo
um instrumento fundamental para a gestão integrada e eficiente do
SIN~\cite{ons:23}.

\section{Algoritmo SDDP}
O algoritmo SDDP (\textit{Stochastic Dual Dynamic Programming}),
proposto por~\cite{MvLm:91}, decompõe o problema de planejamento
hidrotérmico em uma malha bidimensional de subproblemas
determinísticos: cada subproblema é indexado por um par
$(\text{estágio}, \text{cenário})$, em que os estágios representam
períodos discretos da operação --- meses, no caso da operação
programada do SIN --- e os cenários representam realizações
estocásticas das variáveis relevantes, principalmente afluências aos
reservatórios. A solução do problema é construída de forma iterativa,
combinando simulações ao longo dos cenários e refinamentos da política
de operação até a convergência.

A Figura~\ref{fig:SDDP_cenarios_estagios} esquematiza a estrutura de
subproblemas resolvida pelo SDDP nas configurações adotadas nos
experimentos desta dissertação: três estágios mensais e $1\,536$
cenários, totalizando $4\,608$ subproblemas. Os cenários são, por
construção, independentes entre si --- cada cenário corresponde a uma
realização estocástica distinta, sem compartilhamento de estado com os
demais. Essa propriedade caracteriza o SDDP como uma aplicação do tipo
\textit{bag-of-tasks}~\cite{Cirne:2003}: um conjunto de tarefas
mutuamente independentes que podem ser distribuídas entre processadores
sem coordenação durante a execução. Os cenários são particionados entre
os \textit{ranks} MPI e resolvidos em paralelo, de modo que o tempo
total da fase de simulação fica limitado pelo cenário mais lento e não
pela soma sequencial de todos. A natureza \textit{bag-of-tasks} do SDDP
é o que torna o algoritmo escalável em sistemas HPC; é também o que
desloca o gargalo da computação para a persistência dos resultados, à
medida que o número de cenários cresce.

Após a resolução, os resultados de cada subproblema --- variáveis
primais, multiplicadores duais e custos --- precisam ser persistidos em
disco, pois alimentam etapas subsequentes do próprio algoritmo e
análises estatísticas posteriores.

A granularidade temporal adotada para os resultados de cada estágio tem
impacto direto sobre o volume de dados produzido por iteração. A
abordagem tradicional do SDDP agrega a operação dentro de cada estágio
em poucos blocos representativos de carga (\textit{patamares}: pesado,
médio e leve), reduzindo cada estágio a um pequeno número de pontos
operativos por cenário. A operação programada moderna, entretanto,
especialmente com a crescente penetração de fontes intermitentes como
solar e eólica, requer resolução horária para capturar adequadamente
horários de ponta, vales e rampas associadas à variabilidade dessas
fontes. Esta dissertação foca exatamente no regime de resultados
horários, em que cada estágio de um cenário gera dados proporcionais ao
número de horas simuladas --- ordem de magnitude maior que o caso
tradicional por blocos. É nesse regime que o gargalo de E/S se torna
efetivamente dominante: o volume agregado escrito por iteração cresce
linearmente com o número de cenários, com o número de estágios e com a
granularidade temporal interna de cada estágio, e a estratégia de
implementação dessas escritas determina diretamente a escalabilidade da
aplicação em ambientes HPC.

\begin{figure}[H]
  \centering
  \includegraphics[width=\textwidth]{SDDP_cenarios_estagios.png}
  \caption{Estrutura de cenários e estágios do SDDP nos experimentos
  desta dissertação. Cada célula corresponde a um subproblema
  $(\text{estágio}, \text{cenário})$. Os cenários são mutuamente
  independentes e podem ser distribuídos entre \textit{ranks} MPI,
  caracterizando uma aplicação \textit{bag-of-tasks}.}
  \label{fig:SDDP_cenarios_estagios}
\end{figure}

\subsection{Arquitetura atual de E/S do SDDP com um único processo escritor}
\label{sec:arq_atual}
Na arquitetura atual do SDDP, apenas um \textit{rank} MPI (o
\textit{rank}~1) é designado para realizar todas as operações de E/S do
sistema. Cada nó de computação, após resolver sua respectiva tarefa, envia
os resultados por meio de um \textit{buffer} de dados utilizando a função
\texttt{MPI\_Send}. Esse \textit{buffer} é então recebido pelo processo
escritor através de uma chamada \texttt{MPI\_Recv} e, posteriormente, os
dados são gravados em um sistema de arquivos distribuído. A
Figura~\ref{fig:SDDP_architecture} ilustra o funcionamento dessa
arquitetura.

\begin{figure}[H]
  \centering
  \includegraphics[width=0.7\textwidth]{SDDP_architecture.png}
  \caption{Algoritmo SDDP com processo escritor.}
  \label{fig:SDDP_architecture}
\end{figure}

Independentemente do número de nós computacionais utilizados, apenas um
processo é responsável pelas operações de E/S. Cada nó utiliza sua banda
efetiva para enviar os \textit{buffers} de dados, enquanto o processo
escritor utiliza sua banda para recebê-los. A primeira hipótese para as
limitações de escalabilidade nessa arquitetura está relacionada ao
adaptador de rede do servidor que executa o processo mestre.

Adicionalmente, o \textit{rank}~0 é responsável pelo gerenciamento e
controle da aplicação, como a distribuição de tarefas para os demais
\textit{ranks} e o envio/recebimento das configurações do sistema. Essas
atividades também utilizam o mesmo adaptador de rede empregado na
recepção dos \textit{buffers} de dados. Outra hipótese é que o processo
escritor opera continuamente em um fluxo sequencial de recepção e escrita
dos dados, o que pode representar um gargalo adicional, sobretudo em
sistemas com alto volume de dados. Por fim, se houver atrasos na recepção
ou na gravação dos \textit{buffers} de dados por parte do processo mestre,
os nós computacionais podem ser obrigados a aguardar, gerando sobrecarga
no envio dos \textit{buffers} e aumentando o custo computacional da
execução paralela.

